\ifdefined\isMainDocument
\else
\documentclass[11pt]{article}
\usepackage{fullpage}
\usepackage{amsmath,amsthm,amsfonts,amssymb,amscd}
\usepackage{bm}
\usepackage{bbm}
\usepackage{mathtools}
\usepackage{mathrsfs}
\usepackage{etoolbox}
\usepackage{nicefrac}
\usepackage{wrapfig}
\usepackage{lineno}
\usepackage{caption}
\usepackage{graphicx}
\graphicspath{}

\usepackage{geometry}
\geometry{top=0.5in, bottom=0.5in, left=0.5in, right=0.5in}
% \geometry{top=1.0in, bottom=1.0in, left=1.0in, right=1.0in}

\usepackage{setspace}
\setstretch{1.5}

% \usepackage{helvet}
% \renewcommand{\familydefault}{\sfdefault}

\usepackage{fancyhdr}
\pagestyle{fancy}
\fancyhf{} % Clear all headers and footers
% \fancyfoot[R]{\thepage} % Set the page number on the right side of the footer
\renewcommand{\headrulewidth}{0pt} % Remove header rule
\setlength{\footskip}{0pt}

\usepackage[backend=bibtex, style=numeric, sorting=none, sortcites=true, maxnames=5]{biblatex}
\DeclareNameAlias{default}{last-first}
\DeclareFieldFormat{pages}{#1}
\DeclareCiteCommand{\cite}
  {\usebibmacro{prenote}}
  {\usebibmacro{citeindex}%
   \printtext[bibhyperref]{\textsuperscript{\textbf{\printfield{labelnumber}}}}}
  {\multicitedelim}
  {\usebibmacro{postnote}}
\renewbibmacro{in:}{}
\renewcommand*{\multicitedelim}{\textsuperscript{,}}
\addbibresource{refs/refs.bib}
\begin{document}
\appendix
\fi

%---------------------------------------------------------------------

\section{Appendix C: Analytical Integration of the Chromosome-Wide Constraint}
\label{app:chromosapp}
\linenumbers

This appendix derives the chromosome-wide average selective interference, $\overline{Q^2}_\text{linked}$, resulting from heritable fitness variance distributed continuously across a genetic map. The derivation proceeds in two stages. First, we formulate the integral under the exact Haldane mapping function—the expression we numerically evaluate to generate the predictions for every chromosome in our dataset. Second, we derive a closed-form analytical expression under the linear approximation ($r \approx \ell$), formally establishing its biological domain of validity and proving that it provides a strict lower bound on the magnitude of selective interference.

\subsection*{C.1\quad Position-specific constraint profile}

Consider a neutral locus located at genetic position $y$ on a chromosome of total map length $M$ Morgans. The cumulative selective interference exerted on this neutral site by a single locus under selection at position $x$ depends on the persistence of their statistical association. The joint probability that this association survives both recombination and the loss of selective variance across one generation is $(1 - r(\ell))Z$, where $\ell = \vert{}x-y\vert{}$ is the pairwise genetic distance.

The cumulative lifetime interference generated by this selected locus is the infinite geometric sum of this survival probability:

\begin{equation}
    \sum_{t=0}^{\infty}\left[(1 - r(\ell))\,Z\right]^t = \frac{1}{1-(1-r(\ell))Z}
\label{eq:geom_sum}
\end{equation}

Assuming selected loci are distributed uniformly across the genetic map (consistent with the infinitesimal model), we obtain the position-specific interference profile $Q^2(y)$ by squaring this cumulative effect and averaging over all possible positions of the selected locus, $x \in [0,M]$:

\begin{equation}
    Q^2(y) = \frac{1}{M}\int_0^M \left(\frac{1}{1-(1-r(|x-y|))Z}\right)^2 dx
\end{equation}

Under the Haldane mapping function, $r(\ell) = \frac{1}{2}(1-e^{-2\ell})$, the per-generation survival probability becomes:

\begin{equation}
    (1-r(\ell))\,Z = \frac{Z}{2}(1+e^{-2\ell})
\end{equation}

To resolve the absolute value in the genetic distance $\ell = \vert{}x-y\vert{}$, we split the spatial integral at the focal position $y$, computing the interference contributed by the chromosome arms to the left and right of the neutral site separately:

\begin{equation}
    Q^2(y) = \frac{1}{M}\left[\int_0^y f_\text{H}(\ell)\,d\ell + \int_0^{M-y} f_\text{H}(\ell)\,d\ell\right]
\end{equation}
where $f_\text{H}(\ell) = \left(1 - \frac{Z}{2}(1+e^{-2\ell})\right)^{-2}$. Because $f_\text{H}(\ell)$ is a transcendental function of the genetic distance $\ell$, the antiderivative cannot be expressed using elementary functions. Consequently, Equation (C.4) is evaluated numerically.

\subsection*{C.2\quad Chromosome-wide average under the Haldane mapping function}

To determine the aggregate interference experienced by neutral variation across the entire chromosome, we average the position-specific profile $Q^2(y)$ over all possible neutral-site positions, assuming a uniform physical density of neutral sites ($1/M$):
\begin{equation}
    \overline{Q^2}_\text{linked} = \frac{1}{M}\int_0^M Q^2(y)\,dy = \frac{1}{M^2}\int_0^M\int_0^M f_\text{H}(|x-y|)\,dx\,dy
\label{eq:double_int}
\end{equation}

We can reduce this double integral over the spatial domain $[0,M]^2$ to a single integral over all pairwise genetic distances $\ell = \vert{}x-y\vert{}$. Because the interference function is symmetric with respect to direction ($f_\text{H}(\vert{}x-y\vert{}) = f_\text{H}(\vert{}y-x\vert{})$), the contributions from loci where $x > y$ equal those where $x < y$:

\begin{equation}
\overline{Q^2}_\text{linked} = \frac{2}{M^2}\int_0^M\int_0^x f_\text{H}(x-y)\,dy\,dx
\end{equation}

Substituting $\ell = x - y$ in the inner integral ($dy = -d\ell$; moving the limits $y \in [0,x]$ to $\ell \in [x,0]$, which reverses the sign):

\begin{equation}
    = \frac{2}{M^2}\int_0^M\int_0^x f_\text{H}(\ell)\,d\ell\,dx
\end{equation}

The integration region is $\{(\ell,x) : 0 \leq \ell \leq x \leq M\}$. Exchanging the order of integration to evaluate the spatial limits before the genetic distance:

\begin{equation}
    \overline{Q^2}_\text{linked} = \frac{2}{M^2}\int_0^M f_\text{H}(\ell)\left(\int_\ell^M dx\right)d\ell = \frac{2}{M^2}\int_0^M (M-\ell)\,f_\text{H}(\ell)\,d\ell
\label{eq:haldane_avg}
\end{equation}

The biological intuition behind the linear weight $(M-\ell)$ is straightforward: on a chromosome of finite length $M$, the number of locus pairs separated by a given genetic distance $\ell$ decreases linearly as $\ell$ increases. As the genetic distance approaches the total length of the chromosome ($\ell \to M$), the number of physically possible pairs vanishes. This corresponds to Equation (10) in the main text.

\medskip
\noindent\textit{Verification of Biological Boundaries.} The integrity of the model requires that it recovers known biological limits. Using the identity $\frac{2}{M^2}\int_0^M(M-\ell)\,d\ell = 1$, we evaluate the boundaries:

\medskip
\noindent\textit{Complete-linkage limit} ($M \to 0$): For a chromosome with zero recombination (e.g., the avian W or mammalian Y chromosome), all pairwise genetic distances vanish ($\ell \to 0$). The survival function evaluates to $f_\text{H}(0) = (1-Z)^{-2}$ uniformly across the chromosome:

\begin{equation}
    \lim_{M\to 0}\overline{Q^2}_\text{linked} = (1-Z)^{-2}
\end{equation}

\noindent\textit{Independent-assortment limit} ($M \to \infty$): For an infinitely long chromosome, the Haldane function saturates at $r \to 1/2$. Consequently, $f_\text{H}(\ell) \to (1-Z/2)^{-2}$ for the vast majority of locus pairs. Because the exponential convergence to this asymptote is rapid, the dominated convergence theorem ensures:

\begin{equation}
    \lim_{M\to\infty}\overline{Q^2}_\text{linked} = \frac{1}{(1-Z/2)^2}
\end{equation}

These boundaries establish that as a chromosome becomes infinitely long, the selective interference precisely decays to the Robertson unlinked baseline derived in Appendix A ($Q^2 = 4$ as $Z \to 1$).

\subsection*{C.3\quad Closed-form approximation under the linear recombination model}

For genetic distances sufficiently short that interference is severe ($\ell \ll 0.5$ Morgans), we can approximate the Haldane function with a linear expansion $r(\ell) \approx \ell$. Under this linear mapping, the per-generation survival probability becomes $(1-\ell)Z$, yielding the integrand:

\begin{equation}
    f_\text{L}(\ell) = \left(\frac{1}{1-Z+\ell Z}\right)^2
\end{equation}

This approximation allows an analytical antiderivative. Using the substitution $u = 1-Z+\ell Z$, $du = Z\,d\ell$:

\begin{equation}
    \int_0^D f_\text{L}(\ell)\,d\ell = \frac{1}{Z}\int_{1-Z}^{1-Z+DZ}\frac{du}{u^2} = \frac{D}{(1-Z)(1-Z+DZ)}
\label{eq:antideriv}
\end{equation}

Applying this to the left and right chromosome arms ($D = y$ and $D = M-y$) yields the closed-form profile for a specific neutral site:

\begin{equation}
    Q^2_\text{L}(y) = \frac{1}{M(1-Z)}\left(\frac{y}{1-Z+yZ} + \frac{M-y}{1-Z+(M-y)Z}\right)
\end{equation}

To determine the chromosome-wide average, we integrate over all neutral-site positions:

\begin{equation}
    \overline{Q^2}_\text{L} = \frac{1}{M^2(1-Z)}\int_0^M\left(\frac{y}{1-Z+yZ} + \frac{M-y}{1-Z+(M-y)Z}\right)dy
\end{equation}

By symmetry, both arms contribute identically across the interval $[0,M]$, simplifying the integral to:
\begin{equation}
    \overline{Q^2}_\text{L} = \frac{2}{M^2(1-Z)}\int_0^M \frac{y}{1-Z+yZ}\,dy
\end{equation}

Using the substitution $u = 1-Z+yZ$, $dy = du/Z$, we evaluate the integral:

\begin{align}
    \int_0^M \frac{y}{1-Z+yZ}\,dy
    &= \frac{1}{Z^2}\int_{1-Z}^{1-Z+MZ}\left(1 - \frac{1-Z}{u}\right)du \\
    &= \frac{1}{Z^2}\left[u-(1-Z)\ln u\right]_{1-Z}^{1-Z+MZ}
\end{align}

Evaluating at the bounds and combining logarithmic terms yields:

\begin{equation}
    = \frac{M}{Z} - \frac{1-Z}{Z^2}\ln\left(1+\frac{MZ}{1-Z}\right)
\end{equation}

Substituting this back into the leading coefficient $2/(M^2(1-Z))$ provides the final closed-form expression for chromosome-wide interference under linear recombination:

\begin{equation}
    \overline{Q^2}_\text{L} = \frac{2}{MZ(1-Z)} - \frac{2}{M^2Z^2}\ln\left(1+\frac{MZ}{1-Z}\right)
\label{eq:closed_form}
\end{equation}

\subsection*{C.4\quad Domain of validity and directional bias of the approximation}

\noindent\textit{Agreement at complete linkage ($M \to 0$).} Setting $\epsilon = MZ/(1-Z)$ and applying a Taylor expansion to the logarithm ($\ln(1+\epsilon) = \epsilon - \epsilon^2/2 + O(\epsilon^3)$) yields:

\begin{align}
    \overline{Q^2}_\text{L} &= \frac{2}{MZ(1-Z)} - \frac{2}{M^2Z^2}\left[\frac{MZ}{1-Z} - \frac{1}{2}\left(\frac{MZ}{1-Z}\right)^2 + O(\epsilon^3)\right] \\
    &= \frac{1}{(1-Z)^2} + O(M)
\end{align}

This confirms that the linear approximation recovers the Haldane expectation in the limit of zero recombination.

\medskip
\noindent\textit{Failure at independent assortment ($M \to \infty$).} For highly recombining chromosomes, as $M \to \infty$, the leading term $2/(MZ(1-Z)) \to 0$ and the logarithmic correction decays even faster, driving $\overline{Q^2}_\text{L} \to 0$. This violates the biological reality established in (C.10): selective interference cannot decay below the unlinked baseline imposed by independent assortment. The linear approximation assumes recombination rates can exceed $0.5$, artificially purging associations that should persist across unlinked chromosomes.

\medskip
\noindent\textit{Proof of conservative bias.} Because the exponential decay of physical linkage slows as distance increases, the true recombination fraction is strictly concave ($e^{-2\ell} > 1-2\ell$ for all $\ell > 0$). Therefore, the exact Haldane mapping is always bounded above by the linear approximation:

\begin{equation}
    r_\text{H}(\ell) = \frac{1}{2}(1-e^{-2\ell}) < \ell = r_\text{L}(\ell)
\end{equation}

Because the linear approximation overestimates the true recombination rate for any non-zero distance, it systematically underestimates the per-generation survival probability of the statistical association between loci. Consequently, it underestimates the interference function pointwise across the entire chromosome ($f_\text{L}(\ell) < f_\text{H}(\ell)$).

Integrating this strictly lower function guarantees that:

\begin{equation}
    \overline{Q^2}_\text{L} < \overline{Q^2}_\text{H} \quad \text{for all } M > 0
\end{equation}

This establishes that the closed-form expression is a lower bound on the true severity of selective interference. When mapping this to the biological prediction, underestimating interference produces a conservative overestimate of the effective population size, or $N_e/N$.

\medskip
\noindent\textit{Application to empirical data.} The linear approximation is robust when the characteristic distance over which interference decays ($d^* \approx 1-Z$) is small compared to 0.5 Morgans. Biologically, this expression is highly reliable for avian microchromosomes ($M \lesssim 0.3–0.5$ Morgans), where nearly all pairwise locus interactions fall securely within the linear regime. However, for macrochromosomes spanning multiple Morgans, the discrepancy becomes substantial. Therefore, all primary predictions in the main text are derived via numerical integration of the exact Haldane formulation (Equation C.8), utilizing this closed-form expression strictly as an analytical cross-check for localized, short-range behavior.

\ifdefined\isMainDocument
\else
    \end{document}
\fi